\section{Finite relations and a calculus of conditional permutations}
\label{sec:lifting}

The input to this section is the polygon Boolean algebra, its translation
action, and the identity $A_m=[F_m,F_m]$ proved in
Theorem~\ref{thm:simplicity}.  We construct a finite presentation mapping
onto $A_m$ and establish the algebraic rules used to propagate its
relations.  At this stage its kernel is not known to be central.  The
point is to identify precisely which local identities suffice for the
geometric argument in Section~\ref{sec:propagation} and the transport
argument in Section~\ref{sec:transport}.

Throughout the section an \emph{alphabet} is a subset of the $m$ track
indices.  A conditional permutation on a Boolean set $U\subseteq X$
permutes the tracks at the points of $U$ and fixes their complement.
We call $U$ a \emph{test}.  Coordinates of different tracks are aligned
unless offsets are explicitly specified.  We use
$[g,h]=ghg^{-1}h^{-1}$ and write
\[
 B_m=\{d=(d_1,\ldots,d_m)\in\Gamma^m:\textstyle\sum_i d_i=0\}.
\]
Its elements are the balanced translations.  Since
$\OO/\ZZ$ is infinite cyclic, $B_m$ is free abelian of rank $2(m-1)$.

\subsection{Generators and the finite initial choice}

Let $\alpha=\tau\bmod1$.  Choose a sufficiently small square $Q$ about
the origin in a circular coordinate chart, with side levels in $\OO$.
Our initial finite collection $\mathcal U_0$ consists of $X$, the two
coordinate interval tests with endpoints $0,\alpha$, and the two tests
\[
 Q\cap\{y-\lambda x>0\},\qquad
 Q\cap\{x-\lambda y>0\}.
\]
Boundary conventions make no difference in the Boolean algebra.
Complements are generated using the constant permutations, so they
need not be additional basic tests.

For $U\in\mathcal U_0$ and every even permutation $\sigma$ supported
on at most five tracks, include a generator $c_{U,\sigma}$.  Include
also the balanced translations
\begin{equation}\label{eq:lifting-generators}
 t_{i,r}=t(\varepsilon_r e_i-\varepsilon_r e_m),
 \qquad i<m,\quad r=1,2,
 \qquad
 \varepsilon_1=(\alpha,0),\quad\varepsilon_2=(0,\alpha).
\end{equation}
Here $e_i$ designates a track, and the notation $t(d)$ denotes the
corresponding piecewise translation.  This is a finite list, denoted
by $S_m$.  The conditional permutations belong to $A_m$.  The balanced
translations do as well: if $i,j,k,h$ are distinct, $v\in\Gamma$, and
$\sigma=(i\ j\ k)$, then in $F_m$
\begin{equation}\label{eq:translation-commutator}
 [\sigma,t(v e_i-v e_h)]=t(v e_j-v e_i).
\end{equation}
The right side is a commutator in $F_m$, and these differences generate
$B_m$.

\begin{lemma}\label{lem:lifting-generation}
For sufficiently large finite $m$, the list $S_m$ generates $A_m$.
Every polygonal test on a proper alphabet of at least five tracks
is generated, in projection to $A_m$, using conditional permutations
on translated basic tests on that alphabet and balanced translations.
\end{lemma}

\begin{proof}
First consider the Boolean algebra.  Translate the coordinate interval
test by $i\alpha$, with $i$ in an integer interval.  The endpoints of
these tests form a connected path in the orbit-index graph: each test
joins $i\alpha$ to $(i+1)\alpha$.  Two different complementary arcs
of the resulting endpoint set have different test assignments.  Indeed,
an arc between points with identical assignments contains either both
endpoints of every such edge or neither; connectedness then says it
contains all endpoints or none.  The two points therefore belong to
the same complementary arc.  Enlarging the index interval consequently
produces every circular interval whose endpoints lie in $\OO/\ZZ$.
The two coordinate families generate the rectangular Boolean algebra.

Every allowed inclined line is a translate of its line through the
origin.  The allowable tangent translations are dense along the line
by Lemma~\ref{lem:arithmetic}.  Translated copies of $Q$ therefore
cover each relevant compact segment of that line.  In each of finitely
many such charts the translated basic sign test gives its side decision;
rectangular tests clip the decision to the chart.  A finite cover of
the finitely many edges in a polygonal description proves that all
polygonal tests are Boolean combinations of translates of
$\mathcal U_0$.

Fix an alphabet $I$ with $|I|\ge5$.  The group $\Alt(I)$ is perfect.
For completeness, it is generated by $3$-cycles, and the commutator of
two $3$-cycles meeting in one letter is a $3$-cycle.  Conjugating this
identity gives all the generators.  If copies of this group on $U$
and $V$ are available in $A_m$, then
\[
 [g\text{ on }U,h\text{ on }V]=[g,h]\text{ on }U\cap V.
\]
Perfection gives every conditional permutation on $U\cap V$.
The product of the constant $g$ and the inverse of $g$ on $U$ gives
$g$ on $U^c$.  This constructs all Boolean combinations in projection.
A common translation on $I$ can be balanced on one track outside $I$.
Hence it constructs all translated tests as well.

An alternating permutation of slots is a product of permutations of
three slots.  Subdivide their common parameter piece if necessary and
add two spare slots, so that each such permutation lies in an
alternating group on five slots.  When their tracks are distinct, a
balanced translation aligns the five offsets; the conditional group
on the parameter piece is already available.  If tracks repeat, move
each slot to a private track using a conditional $3$-cycle with two
private auxiliary tracks.  These operations do not interfere on a
repeated original track because the slots there are disjoint.  After
this routing the preceding distinct-track construction applies, and
conjugating back gives the original slot permutation.  Five slots
require only finitely many private tracks.  Processing each generator
and each subdivision piece separately proves the assertion for one
fixed sufficiently large $m$.
\end{proof}

We describe carefully how finite relations will be chosen.  If
$\mathcal V$ is any fixed finite list of polygonal tests and $I$ is
a fixed proper alphabet, its pointwise conditional group is finite:
it is
\begin{equation}\label{eq:pointwise-table}
 C_{\mathcal V,I}=\prod_{P\in\operatorname{At}(\mathcal V)}\Alt(I),
\end{equation}
where $\operatorname{At}(\mathcal V)$ consists of the nonzero Boolean
atoms.  By Lemma~\ref{lem:lifting-generation}, choose a word $w_g$
representing each $g$ in this finite group.  Choose it as a product
of conditional permutations on translated basic tests on $I$;
the translations in these expressions are balanced outside $I$.
Expanding gives a finite word on the original generators.
Impose the finitely many true relations
\begin{equation}\label{eq:table-relations}
 w_1=1,\qquad w_gw_h=w_{gh},
\end{equation}
and identify each specified input conditional permutation with its
chosen representative.  The resulting copy of the table is exact:
the table relations give a homomorphism, and projection to $A_m$ is
its left inverse.  Several chosen descriptions of the same bounded
configuration can also be identified by finitely many true relations.

Here and below \emph{bounded data} means data belonging to one fixed
finite list before any propagation or arbitrary word is considered.
There is no assertion that tests of arbitrarily large labels already
have their tables.  The initial list of required relations comprises:
\begin{enumerate}[label=(\roman*)]
\item the joint conditional alternating-group tables for
      $\mathcal U_0$, their compatibility on subalphabets,
      disjoint-alphabet commutation, and their covariance under
      constant alternating track permutations;
\item commutation of the $t_{i,r}$, giving the additive homomorphism
      $t:B_m\to\widehat G$, and constant-permutation covariance
      $c_{X,\sigma}t(d)c_{X,\sigma}^{-1}=t(\sigma d)$, checked on
      the finite bases;
\item commutation of a basic conditional copy with translations
      vanishing on its alphabet, checked on a basis of those
      translations, and invariance of a coordinate test under
      translations in the other coordinate;
\item the finite tables, with their specified overlap identities,
      for the geometric configurations described below, on all
      alphabets needed for the support and conditional-move arguments;
\item the identities \eqref{eq:translation-commutator} for the basis
      differences in \eqref{eq:lifting-generators}.
\end{enumerate}
All these relations hold in $A_m$.

The finite geometric menu in (iv), including its initial coordinate
window, is specified in Section~\ref{sec:propagation}. That section
orders the choices and proves that common translates of this fixed
menu suffice at every later stage. Each of its regions has a finite
expression in translated basic tests by
Lemma~\ref{lem:lifting-generation}, so the table construction above
applies. The menu is fixed before any propagation or later word is
considered; the algebraic rules below apply to any such finite choice.

For the resulting fixed $m$ and menu, every relation above is a finite
word on $S_m$. If $L$ is their maximum length, one possible single
presentation is
\begin{equation}\label{eq:finite-presented-cover}
 \widehat G=\langle S_m\mid R_L\rangle,
 \qquad
 R_L=\{w:|w|\le L,\ w=1\text{ in }A_m\}.
\end{equation}
The set $R_L$ is finite and consists of true relations.  Its use is
existential; no decision procedure for membership in $R_L$ is needed.
Lemma~\ref{lem:lifting-generation} supplies the epimorphism
$\pi:\widehat G\twoheadrightarrow A_m$.

\subsection{Perfect covers and the meaning of a central law}

The exact initial tables will eventually yield relations only up to
central elements.  Perfect covers make their conditional factors
unambiguous even in this situation.

\begin{lemma}\label{lem:perfect-covers}
Every perfect group $E$ has a canonical perfect central extension
$E^\ast\twoheadrightarrow E$.  It is functorial in $E$.  A map from
$E$ into the quotient of any central extension has a unique lift from
$E^\ast$.  Moreover:
\begin{enumerate}[label=(\alph*)]
\item two homomorphisms from a perfect group into a central extension
      that induce the same quotient map are equal;
\item if two perfect subgroups commute modulo a central subgroup,
      they commute exactly;
\item a perfect group $H$ has centerless quotient $H/Z(H)$.
\end{enumerate}
\end{lemma}

\begin{proof}
Take a free presentation $E=P/R$ and put
\[
 E^\ast=[P,P]/[P,R].
\]
Its map to $E$ is surjective because $E$ is perfect, and its kernel
is central.  In $P/[P,R]$, the image of $R$ is central and every
element is a product of an element of the derived group and an
element of that image.  Taking commutators twice therefore shows
that the derived group is perfect.  This proves perfection of
$E^\ast$.  Given a central lifting problem, choose free lifts of
the generators of $P$.  They kill $[P,R]$, so their restriction to
$[P,P]$ descends to the required homomorphism.  Changing the free
lifts multiplies them by central elements and hence does not change
this restriction.  Applying (a) to two lifts from the perfect group
$E^\ast$ proves uniqueness.  This gives independence of presentation
and functoriality.

For (a), the pointwise ratio of the two homomorphisms has central
values and is consequently a homomorphism to an abelian group; it
vanishes on a perfect group.  For (b), when the relevant commutators
are central, commutation with a fixed element is a homomorphism on
the other subgroup.  Perfection makes it trivial.  Finally, if the
image of $h\in H$ is central in $H/Z(H)$, then
$x\mapsto[h,x]$ has central values and is a homomorphism.  It is
trivial because $H$ is perfect, so $h\in Z(H)$, proving (c).
\end{proof}

This is the universal central-extension construction; compare
\cite[Construction~6.9.3 and Theorem~6.9.5]{Weibel1994}.
We retain the construction here
because its uniqueness is used repeatedly.

Fix tests $U_1,\ldots,U_r$, an alphabet $I$ with $|I|\ge5$, and a
set $\Omega\subseteq\{0,1\}^r$ of allowed assignments.  Initially
$\Omega$ may contain assignments not realized in $X$.  It must
contain every actual assignment, and we retain any exclusions
already established for a subfamily.  The \emph{assignment group} is
\[
 C_{\Omega,I}=\prod_{\omega\in\Omega}\Alt(I).
\]
A permutation conditional on $U_j$ evaluates as that permutation
in the factors with $\omega_j=1$ and as the identity elsewhere.
The constant group evaluates diagonally.  These groups generate
$C_{\Omega,I}$: commutators construct intersections, complements
follow using the diagonal, and perfection gives every factor.

Let $H$ be the subgroup of $\widehat G$ generated by the specified
conditional copies.  They may be exact copies of $\Alt(I)$ or
images of $\Alt(I)^\ast$ already obtained.  We say they have the
\emph{central law for $\Omega$} if every word that evaluates to
$1$ in $C_{\Omega,I}$ belongs to $Z(H)$.  When star groups are
used, evaluation means evaluation of their underlying alternating
permutations.  Equivalently, the specified generator maps induce
a surjection
\begin{equation}\label{eq:central-law-map}
 C_{\Omega,I}\longrightarrow H/Z(H).
\end{equation}
This definition does not presume a map from $H$ to the formal
assignment group.  In particular it remains meaningful before
unrealizable assignments have been removed.

Applying Lemma~\ref{lem:perfect-covers} to
\eqref{eq:central-law-map} gives a canonical representation
\begin{equation}\label{eq:sector-representation}
 \rho_{P,I}:\Alt(I)^\ast\longrightarrow H\subseteq\widehat G
\end{equation}
for every sector $P$, and for every union of sectors.  A formal
sector is designated by its assignment, even if its geometric
intersection is empty; its representation need not yet vanish.
When the assignments are actual, $P$ denotes the corresponding
Boolean region.

The images for disjoint sectors commute, and if $P$ is a disjoint
union of sectors $P_a$, then
\begin{equation}\label{eq:canonical-splitting}
 \rho_{P,I}(s)=\prod_a\rho_{P_a,I}(s),
 \qquad s\in\Alt(I)^\ast.
\end{equation}
Indeed the factors commute by Lemma~\ref{lem:perfect-covers}(b),
and the two homomorphisms in \eqref{eq:canonical-splitting} have
the same map to $H/Z(H)$, so (a) applies.  These images generate
$H$: their product $H_0$ maps onto $H/Z(H)$, whence
$H=H_0Z(H)$, and perfection of $H$ gives $H=H_0$.  The same
uniqueness proves compatibility with refinement and restriction to a
subalphabet within this lawful family, and with any previously specified
perfect copy in that family.
An excluded factor is trivial because
its perfect image is central.

Later we shall enlarge a lawful family and exclude some assignments
that were previously allowed. The centers of the two generated groups
need not agree, so this comparison requires a separate argument.

\begin{lemma}[Compatibility under enlargement]\label{lem:law-enlargement}
Fix an alphabet $I$ with $|I|\ge5$. Suppose a finite test family
$\mathcal T_0$ is contained in a finite family $\mathcal T_1$, with
literally the same specified input copies on the old tests, including
the constant copy. Let $H_0\le K$ be their generated subgroups.
Suppose they have central laws for assignment sets $\Omega_0$ and
$\Omega_1$, respectively, and restriction of assignments gives a map
$r:\Omega_1\to\Omega_0$. Then their canonical representations satisfy
\[
 \rho^{\,0}_{\omega,I}(s)
 =\prod_{\substack{\omega'\in\Omega_1\\r(\omega')=\omega}}
       \rho^{\,1}_{\omega',I}(s),
 \qquad \omega\in\Omega_0,\quad s\in\Alt(I)^\ast.
\]
The factors on the right commute. An old assignment with no extension
has trivial original representation.
\end{lemma}

\begin{proof}
Put $C_0=C_{\Omega_0,I}$ and $C_1=C_{\Omega_1,I}$. Restriction induces
$d:C_0\to C_1$: an old factor is repeated on all assignments extending
it, and is killed if there are none. Write
$\Lambda:C_0^\ast\to H_0$ for the old canonical perfect-cover map,
$p:C_0^\ast\to C_0$ for its cover projection, and $q:K\to K/Z(K)$.
The new law gives $\theta:C_1\to K/Z(K)$.

The homomorphisms $q\Lambda$ and $\theta dp$ agree on the canonical
lifts of every defining old input copy: those inputs are the same
elements of $H_0$ and $K$. These lifts generate $C_0^\ast$. Indeed,
their generated subgroup $M$ surjects onto $C_0$, because the defining
conditional groups generate the assignment group. Thus
$C_0^\ast=M\ker p$; centrality of $\ker p$ and perfection give
$C_0^\ast=[M,M]\le M$. Consequently
\[
 q\Lambda=\theta dp.
\]
For an old assignment with no extension, this equality puts its
original perfect image in $Z(K)$, so that image is trivial. Otherwise
it identifies the quotient maps of its old representation and the
product of its extending new representations. The new factors commute,
so their product is a homomorphism from $\Alt(I)^\ast$. Uniqueness of
lifts from a perfect source into $K\to K/Z(K)$ identifies the two
homomorphisms exactly. No inclusion $Z(H_0)\subseteq Z(K)$ is needed.
\end{proof}

We will also use that $\Alt(I)^\ast$ is generated by the images
of $\Alt(J)^\ast$ for five-element $J\subseteq I$.  Those images
generate a subgroup surjecting onto $\Alt(I)$.  It is invariant
under conjugation, by functoriality and uniqueness; its product
with the central kernel is all of $\Alt(I)^\ast$.  Perfection
then gives the assertion.  Thus every centrality or covariance
test for a canonical representation can be made one five-track
copy at a time.

\subsection{The four algebraic rules}

The following rules separate finite track bookkeeping from the
geometric propagation.  In particular, the length of a Boolean
expression never measures the number of tracks it requires.

\begin{lemma}[Small-support presentations]\label{lem:small-support}
There is an absolute integer $q$ with the following property.
Suppose conditional perfect copies are specified coherently on
the subalphabets of a finite track set, with a common assignment
model $\Omega$.  If the central law holds on every alphabet of
at most $q$ tracks, then it holds on the entire track set.
It suffices to take $q=15$ for this implication.  Additional
tracks used to prove its hypotheses are a separate fixed reserve.
The conclusion concerns the subgroup generated by the indicated
test family, not translations or other predicates outside that
family.
\end{lemma}

\begin{proof}
We give the support count, since it prevents reserves from growing
under repeated refinement.  Present $\Sym(I)$ by transpositions
$t_{ab}$, imposing their involution relations, disjoint
commutation, and the conjugation relations
\[
 t_{ab}t_{bc}t_{ab}=t_{ac}\quad(a,b,c\text{ distinct}).
\]
These relations present the symmetric group: adjacent
transpositions satisfy the usual braid and distant commutation
relations, whose normal form moves the final letter to its
prescribed position and then proceeds inductively; conversely
the displayed relations express every transposition in adjacent
ones.  Each defining relation involves at most four letters.

Fix a five-letter reserve $E\subseteq I$ and a transposition
$a$ on two of its letters.  Reidemeister--Schreier rewriting for
the parity subgroup, with coset representatives $1,a$, gives
generators $a t_{bc}$ and their inverses.  Each is an even
permutation on at most four letters.  To see the support bound
directly, insert the running representative, $1$ or $a$, between
successive letters of a symmetric relator.  The rewritten
relations use only the original at most four letters and the
two letters of $a$.  Explicitly, using the equivalent symmetric
relations $x_tx_sx_t=x_{tst}$ indexed by transpositions, the
rewritten presentation with $u_t=at$ is
\[
 u_a=1,\qquad
 u_t^{-1}u_su_t^{-1}u_{tst}=1,\qquad
 u_tu_s^{-1}u_tu_{tst}^{-1}=1
 \quad(s,t\text{ transpositions}).
\]
The involution relations rewrite to cancellations.
This also proves the subgroup presentation: insert the running
representatives in a product of conjugates of symmetric relators
representing a null word.  No letter outside those in the relator
and $a$ is introduced.

For each $\omega\in\Omega$, realize its conditional copy of
$a t_{bc}$ by a word in the test copies on
$E\cup\{b,c\}$, an alphabet of size at most seven.  Such a
word exists by the intersection-and-complement argument preceding
\eqref{eq:central-law-map}; its length can depend on $\Omega$.
Use the rewritten alternating presentation on each assignment,
and cross commutators between distinct assignments.  Every
relator so lifted involves at most the five reserve letters
and four other letters, hence at most nine tracks.

An original generator on five tracks can be compared with the
product of its assignment generators on the union of those
tracks with $E$, of size at most ten.  Include all these
comparisons and the internal relations of the input perfect
copies.  The resulting relations present the assignment group:
after the comparisons every input is expressed by assignment
generators, and the preceding product presentation then applies.
Equivalently, the kernel of evaluation from the free group on
the specified input symbols is normally generated by null
words supported on at most ten tracks.  This argument also
applies to star inputs: their internal words evaluate through
the underlying alternating permutations, and remain supported
on the same five tracks.

Let $r$ be any one of these small null words and let $g$ belong
to an original five-track perfect copy.  Their combined
alphabet has at most fifteen tracks.  The central law there
says $[r,g]=1$.  Thus every normal generator of the evaluation
kernel is central in the full subgroup $H$.  All its conjugates
are therefore central, giving the global central law.  For
$|I|<5$ one embeds in a five-letter alphabet; for
$5\le|I|\le15$ the conclusion is already among the hypotheses.
The proof never requires more tracks for longer atom words or
for more assignments.
\end{proof}

This argument explains a convention used henceforth.  Laws are
proved simultaneously for small alphabets, including one extra
five-track input whenever centrality is tested, and are then
extended by Lemma~\ref{lem:small-support}.  A central error in
one small subgroup has not been declared central in a larger
one without this additional test.

\begin{lemma}[Translated primitives]\label{lem:translated-copies}
Let $U$ be a basic test and let $I$ be a proper alphabet.  For
$u\in\Gamma$, choose $d\in B_m$ with $d_i=u$ for $i\in I$.
Then
\begin{equation}\label{eq:translated-copy}
 \rho_{U+u,I}(s)
   =t(d)\rho_{U,I}(s)t(d)^{-1}
\end{equation}
is independent of the balancing values off $I$.  The same
assertion holds for words and canonical perfect factors obtained
from a lawful family on $I$.  It has the following consequences.
\begin{enumerate}[label=(\alph*)]
\item Conditional copies on disjoint small alphabets commute
      exactly, even when their tests have different translation
      offsets.  Products of such translated input copies inherit
      this property.
\item Translating all tests in a fixed finite table by a common
      $u$ preserves its exact law on a proper alphabet.  Thus
      finite relations for bounded offsets give laws for those
      offsets relative to any common shift.
\item A translation whose coordinates vanish on $I$ commutes
      with all such aligned words and perfect factors on $I$.
      Coordinate predicates are unchanged by a shift in the
      other coordinate.
\end{enumerate}
\end{lemma}

\begin{proof}
Two choices of $d$ differ by a balanced translation zero on $I$.
The relations (ii)--(iii), first on a basis and then by additivity,
make this difference commute with the basic copy.  Since all
translations commute, the same is true of any translate, then
of every word on those translated copies.  Canonical factors
are represented by such words, so the assertion includes them.

For disjoint alphabets $I,J$, choose one balanced vector that
equals $u$ on $I$ and $v$ on $J$, compensating on a track
outside their union.  Both translated groups are conjugates
by this single vector of their unshifted groups, which commute
by the finite initial tables.  Balance independence identifies
them with any other chosen descriptions.  Conjugating a word
or a table proves (b); uniqueness of perfect lifts identifies
the conjugated canonical factors with those for the translated
family.  The preceding commutation proves (c).
The last statement in (c) follows from its imposed
generating-shift instances and additivity.
\end{proof}

An \emph{aligned word on $I$} will mean a word in these translated
conditional copies whose indicated track supports are contained
in $I$.  The balancing words in the original finite generators
may mention other tracks.  Lemma~\ref{lem:translated-copies}
shows that this does not enlarge the support used in the
disjoint-commutation rule.  A newly obtained canonical factor
on $I$ belongs to the subgroup of aligned words on $I$.

For example, an arbitrary translate of an $x$-coordinate basic
test and an arbitrary translate of a $y$-coordinate basic test
have their joint law: translate the fixed two-test table by the
vector having the prescribed $x$ and $y$ coordinates, and use
transverse invariance.  This observation starts the rectangular
propagation without requiring unbounded initial tables.

We next formalize containment.  Let $J$ be a test with specified
conditional representations on the relevant alphabets.
A family of perfect fragments
$\rho_{P,I}$ is \emph{controlled inside $J$} if conjugation by
a conditional alternating permutation on $J$ agrees on those
fragments with conjugation by the corresponding constant track
permutation.  The definition is tested on small alphabets and
their common enlargements.  Constant conjugation reindexes a
fragment's alphabet; this equivariance follows from the initial
relations and the uniqueness in Lemma~\ref{lem:perfect-covers}.
The label $P$ can still be a formal sector.  Thus control is an
identity about representations, not just a containment after
applying $\pi$.

\begin{lemma}[Control by containers]\label{lem:control}
Assume the relevant small-alphabet laws and the exact
disjoint-alphabet commutation of
Lemma~\ref{lem:translated-copies}.  Assume also that each
individual comparison leaves a fresh five-track bank and two
parity-correction letters outside the alphabets involved,
and one further track for balancing translations.
\begin{enumerate}[label=(\alph*)]
\item If the allowed assignments of a central law imply
      $P\subseteq J$, they give control inside $J$; disjointness
      in that assignment model gives commuting perfect
      representations.  In particular, geometric inclusion
      and disjointness suffice in an actual central law.
\item If a fragment is controlled inside $J$ and an aligned
      word $u$ centralizes the conditional perfect copies on
      $J$ on the enlarged alphabets used below, then $u$
      centralizes that fragment.
\item Control is transitive.  More generally, if two aligned
      words have the same conjugation action on the conditional
      copies on $J$, they have the same action on every
      fragment controlled inside $J$.
\item If $J,K$ have a joint central law and a fragment is
      controlled inside each, it is controlled inside
      $J\cap K$.  Fragments controlled on disjoint sides of
      a lawful decomposition commute, even if those fragments
      do not yet have a joint law with one another.
\end{enumerate}
\end{lemma}

\begin{proof}
Part (a) follows from the commuting sector factors and
\eqref{eq:canonical-splitting}.  A formal predicate side over
an actual rectangular cell satisfies this hypothesis for that
cell and its known rectangular containers: those implications
already hold in the retained rectangular coordinates of the
assignment model.  Geometric emptiness of the formal predicate
side alone is not sufficient.  For (b), test one perfect
copy on an alphabet $I$ of five tracks.  Choose a fresh alphabet
$I'$ disjoint from the track support of $u$ and an even
permutation $\sigma$ carrying $I$ onto $I'$.  If necessary,
two further unused letters correct its parity.  Let $v$ be a
lift of this permutation conditional on $J$, on the union
of the relevant alphabets.  By hypothesis $[u,v]=1$.
Control gives
\[
 v\rho_{P,I}(s)v^{-1}
   =\rho_{P,I'}(\sigma_*s),
\]
where $\sigma_*$ is induced reindexing of the star group.
The right side commutes with $u$ by disjoint-alphabet
commutation.  Therefore
\begin{equation}\label{eq:conditional-move}
 [u,\rho_{P,I}(s)]
 =v^{-1}[u,v\rho_{P,I}(s)v^{-1}]v=1.
\end{equation}
Only this one fragment is moved at a time.

Apply (b) to the ratio of two words to obtain the comparison
assertion in (c).  To obtain transitivity, suppose $P$ is
controlled inside $J$ and $J$ inside $K$.  The ratio of a
$K$-conditional permutation and its constant counterpart
centralizes the $J$ copies; hence it centralizes $P$ by (b).

For intersection control, write the four commuting factors
of the joint $J,K$ law as
\[
 A=J\cap K,\quad B=J\cap K^c,\quad
 C=J^c\cap K,\quad D=J^c\cap K^c.
\]
Use the same star source on an alphabet large enough for the
test under consideration, and abbreviate its representations
by $a,b,c,d$.  Control inside $J$ says that every
$c(s)d(s)$ centralizes the fragment; control inside $K$ says
that every $b(t)d(t)$ does.  Taking commutators gives
\[
 [c(s)d(s),b(t)d(t)]=[d(s),d(t)].
\]
The image of $d$ is perfect, so it centralizes the fragment.
It follows that $b$ and $c$ do too.  The $a$ factor therefore
has the same action as $abcd$, the constant permutation.
This proves control inside $J\cap K$ without a joint law
involving the fragment.

Finally, suppose one fragment is controlled on a side $J$
and another on $J^c$.  Move the first to fresh tracks by a
$J$-conditional permutation.  This fixes the second by its
own control and the $J,J^c$ decomposition.  The moved
fragments commute by disjoint alphabets; conjugation back
proves the assertion.  Iterating handles disjoint unions
of sectors.
\end{proof}

Here is the container comparison used later for overlapping charts.
Suppose $U,V,J$ have a joint actual central law and
$U\cap J=V\cap J$. For a fixed $s\in\Alt(I)^\ast$, the element
$u=\rho_{V,I}(s)^{-1}\rho_{U,I}(s)$ centralizes the conditional
perfect copies on $J$, by their commuting sector decomposition on the
required enlarged alphabets. If a fragment $\rho_{P,I'}$ is controlled
inside $J$, the lemma gives
\[
 \operatorname{Ad}(\rho_{U,I}(s))\big|_{\rho_{P,I'}(\Alt(I')^\ast)}
 =\operatorname{Ad}(\rho_{V,I}(s))\big|_{\rho_{P,I'}(\Alt(I')^\ast)}.
\]
The fragment need not have a joint law with $U,V$ and may still be
formal. The proof moves its five-track factors within $J$ to unused
tracks, where disjoint-alphabet commutation is available, and then
moves them back. Thus the known agreement on $J$ yields an equality of
actions on the fragment, beyond an equality of projected supports.

\begin{lemma}[Pair laws and simultaneous refinement]
\label{lem:pair-laws}
Suppose a finite collection of partitions has a central law
for each individual partition and each pair, coherently on
subalphabets, and conditional
copies on disjoint alphabets commute exactly.  Assume the
spare tracks specified in Lemma~\ref{lem:control} are available
for the individual comparisons.  Then the whole
collection has the central law for formal assignments, with
all exclusions already known in any lawful subfamily retained.
Its canonical representations refine all the pairwise
decompositions simultaneously.
\end{lemma}

\begin{proof}
It suffices first to work on an alphabet with fixed spare
tracks and then apply Lemma~\ref{lem:small-support}.  Let
$H$ be generated by the conditional perfect groups.  Choose
a union $J$ of atoms of one partition.  Each input copy has
a separate joint law with $J,J^c$ and therefore an exact
decomposition into its $J$ and $J^c$ factors.  Let $H_J$
and $H_{J^c}$ be generated by all those respective factors.
They are perfect and generate $H$.

They commute.  To check this, take a five-track factor from
each side.  A permutation conditional on $J$ moves the first
factor to fresh tracks, by its own pair law, and fixes the
second, by the latter's pair law.  Their resulting disjoint
alphabets commute.  This calculation uses neither a triple
law nor the desired simultaneous refinement.  Thus
\[
 H=H_JH_{J^c},\qquad [H_J,H_{J^c}]=1.
\]
The images in $H/Z(H)$ form a direct product: an element in
their intersection commutes with both images and hence is
central in $H/Z(H)$, whose center is trivial by
Lemma~\ref{lem:perfect-covers}(c).

We recall why different direct decompositions of a centerless
group refine one another.  Write $Q=A\times B=C\times D$.
For $a\in A$, its $C$ and $D$ components separately commute
with $B$, since they do so after projecting the equality
$[a,B]=1$ into the two factors.  The centralizer of $B$ in
$A\times B$ is $A$, because $B$ is centerless.  Thus each
component of $a$ belongs to $A$.  The same applies to $B$,
so the two decompositions refine into their four
intersections.  Repeat this argument for the finitely many
partition decompositions.

In every resulting factor $Q_\omega$ of $H/Z(H)$ an input
conditional permutation is the constant permutation if its
assigned predicate is true, and the identity otherwise.
The factor is generated by these images, so it is a quotient
of the constant $\Alt(I)$.  For star inputs their central
kernels vanish in this factor: their pair law with each
input group makes those kernels central in $H$.
The diagonal constant source is an exact alternating group
by the initial tables.  Consequently the simultaneous
decomposition supplies the surjection
$C_{\Omega,I}\twoheadrightarrow H/Z(H)$, proving the formal
central law.

We verify separately that earlier exclusions survive.
Suppose the restriction of $\omega$ to a subfamily
$\mathcal T$ is already forbidden.  The usual commutator
and complement words select that assignment in the formal
model of $\mathcal T$, with any prescribed alternating
permutation as value there.  In the already known law for
$\mathcal T$ these words evaluate to identity, so their
images are central in $H_{\mathcal T}$.  In $Q_\omega$,
however, they give every image of the constant alternating
group.  The image of $H_{\mathcal T}$ is the whole of
$Q_\omega$, since it includes that constant group.
It follows that $Q_\omega$ is abelian.  It is also perfect,
as a quotient of $\Alt(I)$, and is therefore trivial.
This removes the forbidden assignment without assuming
that a center of a subfamily is central in all of $H$.

Finally, each new assignment restricts to an allowed assignment of
every old lawful subfamily, by the exclusions just proved. The input
copies in these subfamilies are the same copies generating $H$.
Lemma~\ref{lem:law-enlargement} therefore identifies their canonical
perfect factors with the corresponding products of new factors,
including their unions. Applying
Lemma~\ref{lem:small-support} proves the asserted law on
the full alphabet.
\end{proof}

\subsection{Patching actions on a rectangular partition}

The remaining algebraic issue is local verification of a
relation when different predicates have only separate
decompositions over the same rectangles.  The subgroup
defined next is deliberately formed before assuming any
joint law among those predicates.

Let $\mathcal P$ be a finite rectangular partition, with
its actual central law.  Suppose each input predicate
$U_j$ separately has a joint law with $\mathcal P$.
The latter laws may allow both formal choices of $U_j$
over a cell.  For $P\in\mathcal P$, define $H_P$ to be
the subgroup generated by all restricted perfect copies
over $P$: the $U_j\cap P$ and $U_j^c\cap P$ factors for
every $j$, including their small subalphabet copies.
Each of these factors is controlled inside the genuine
rectangle $P$, even when its predicate side is formal.
Write $H$ for the subgroup generated by the original input copies.

\begin{lemma}[Invariant cell groups and patching]
\label{lem:patching}
In the preceding situation, with the spare tracks of
Lemma~\ref{lem:control} available for each individual comparison,
the groups $H_P$ commute for
distinct cells, generate a subgroup containing $H$, and are
preserved by every original conditional generator.
To prove that a word $w$ in those generators is central
in $H$, it suffices to prove that it acts trivially on
each $H_P$.  The following verification is sufficient:
for every letter of $w$ choose a replacement whose
conjugation action on $H_P$ is the same; require that the
replacement word acts trivially on $H_P$.

Equality of the actions can be established by a chain
of container comparisons from Lemma~\ref{lem:control}.
In particular, a relator central in a lawful template
group acts trivially on $H_P$ if the template contains
a controlling container $J$ for $H_P$ and its central
law is available on the enlarged alphabets needed for
the conditional move.  This remains true when the
individual $U_j$-sides over $P$ are only formal.
\end{lemma}

\begin{proof}
For different cells, take one generating perfect copy
in each.  A permutation conditional on the first cell
moves that copy to fresh letters and fixes the other
one, using their separate laws with $\mathcal P$.
Disjoint-alphabet commutation and conjugation back show
that $H_P$ and $H_{P'}$ commute.  The exact splittings
of the original copies show that the subgroup generated by the
$H_P$ contains $H$.

Fix an original generator $p(s)$ and one cell $P$.
Its own separate law with $\mathcal P$ gives
\begin{equation}\label{eq:cell-invariance-split}
 p(s)=p_P(s)p_{P^c}(s),\qquad p_P(s)\in H_P.
\end{equation}
For a generating five-track fragment $a_P$ of $H_P$,
move its letters off the support of $p_{P^c}(s)$ by a
permutation conditional on $P$.  This is the ordinary
reindexing on $a_P$, by $a_P$'s own $P$-law, and fixes
$p_{P^c}(s)$, by $p$'s own $P$-law.  Disjoint-alphabet
commutation gives
\begin{equation}\label{eq:cell-complement-commutes}
 [p_{P^c}(s),a_P]=1.
\end{equation}
Thus $p$ acts on $H_P$ by its factor $p_P(s)$, an inner
automorphism of $H_P$.  This proves invariance using
only two separate rectangular laws; no law relating
the predicates of $p$ and $a_P$ was used.

All original conjugations can now be composed as
automorphisms of the same $H_P$.  A replacement with
the same action also preserves $H_P$.  Equality of
their individual actions therefore implies equality
of the actions of the two words.  If the replacement
word acts trivially, then $w$ centralizes $H_P$.
Doing this for every cell proves that $w$ centralizes
the group they generate, and therefore $H$.

For clarity, consider one container comparison.
Let $u$ be the ratio of two proposed replacements,
and suppose their known joint law says that $u$
centralizes every conditional perfect copy on $J$.
The fragment $a_P$ is controlled inside $P$, and
actual rectangular containment and
Lemma~\ref{lem:control} make it controlled inside
$J$.  Equation~\eqref{eq:conditional-move} gives
$[u,a_P]=1$.  Intersections of controlling rectangles
are permitted by the intersection part of that
lemma, so the argument includes clipped containers.

Finally let $r$ be a central relator in a template
containing $J$.  The central-law assertion is applied
on the union of the alphabet of $r$, that of the
particular fragment $a_P$, and a fresh bank for the
move.  In that template group $r$ commutes with the
$J$-conditional mover.  After the move, $a_P$ has
disjoint alphabet from $r$, so they commute exactly;
conjugating back proves $[r,a_P]=1$.  This applies
one generator of $H_P$ at a time.  It does not require
$H_P$ to belong to the template group.  Nor does it
promote centrality in a smaller subgroup without
testing the additional letters.  The proof therefore
also applies to formal wrong-side fragments, whose
actual rectangular control was available from the
start.
\end{proof}

We finish by collecting the finite-reserve implication.
The support presentation uses a fixed reserve of five
letters and one extra five-track test.  A conditional
move compares two such supports and uses a fresh copy
of one of them, with at most two parity-correction
letters.  Translation balancing uses one further
track whenever the indicated alphabet is proper.
The later transport argument compares only a fixed
number of five-slot routings at a time.  Each of these
is a fixed finite operation, so the maximum of their
track requirements, together with the homology
stability threshold, is finite.  Choose $m$ once
larger than this maximum.  A large number of atoms,
a long slide, successive refinements, and a long
word are processed one comparison at a time and do
not change $m$ or the initial relation list.

We have obtained all the algebraic rules needed for
propagation.  To complete the construction it remains
to show that the fixed geometric tables force central
laws for arbitrary polygonal tests; this is the task
of the next section.  The distinction maintained
here will be essential there: a formal sector can
already have exact rectangular control, while its
vanishing is still to be proved from geometric
containment.
